\section{Sesgo, robustez y generalización}
\label{sec:marco_generalizacion}

Las métricas de la Sección~\ref{sec:marco_evaluacion} suelen calcularse sobre un conjunto de test de la misma
fuente que el de entrenamiento, pero en uso real las imágenes llegan de otros equipos, centros y pacientes.
Ese cambio de distribución (\emph{distribution shift}) reduce el desempeño de forma sustancial en problemas de
todo tipo, y la caída persiste aun con los métodos diseñados para mitigarla. Se distinguen el cambio de
dominio, cuando las imágenes provienen de una fuente que el entrenamiento no vio, y el cambio de
subpoblación, cuando la fuente es la misma pero los grupos aparecen en otras
proporciones~\cite{koh2021wilds}.

En dermatología el cambio de distribución más visible es el de modalidad. La fotografía clínica registra la
lesión tal como se ve a simple vista, y la dermatoscopía la examina con aumento e iluminación controladas, con
lo que muestra estructuras que la fotografía no resuelve, como la red pigmentaria y los patrones
vasculares~\cite{dinnes2018dermoscopy,kittler2016standardization}. Los rasgos visuales tienen distribuciones
distintas en uno y otro registro, así que un modelo entrenado con una modalidad y usado con la otra enfrenta
un cambio de dominio.

Aun dentro de una misma modalidad, las condiciones de adquisición varían entre equipos y centros: la
iluminación, la distancia a la lesión, la resolución del sensor y la compresión de la imagen. Un modelo puede
apoyarse en rasgos que esas variaciones alteran, como la textura fina o el color exacto, y perder desempeño
ante imágenes de contenido clínico idéntico. Esa fragilidad se estima perturbando las imágenes de test de
manera controlada (Sección~\ref{sec:marco_eval_robustez})~\cite{hendrycks2019benchmarking}.

Un conjunto puede contener además rasgos que acompañan a una clase sin formar parte de la lesión, como las
marcas de tinta con que se señala una lesión antes de biopsiarla, las etiquetas adhesivas y los bordes oscuros
del dermatoscopio. Si en el conjunto esos rasgos aparecen con más frecuencia en los melanomas, un modelo puede
aprender a predecir melanoma cuando los ve y alcanzar una exactitud alta sin mirar la lesión, un atajo que se
ha documentado en conjuntos públicos de dermatoscopía~\cite{yan2023towards}. Los propios modelos contrastivos
médicos permiten buscar esos artefactos en un conjunto antes de entrenar sobre él~\cite{kim2024transparent}, y
un mapa de Grad-CAM que marca la tinta y no la lesión (Sección~\ref{sec:marco_eval_gradcam}) es el síntoma
visible del atajo.

El tono de piel es un caso de cambio de subpoblación. Cuando las pieles claras están mucho más representadas
que las oscuras en el entrenamiento, el modelo clasifica mejor los fototipos parecidos a los que
vio~\cite{groh2021neural}. Evaluados sobre imágenes con fototipos diversos y confirmación histopatológica,
modelos dermatológicos del estado del arte rindieron bastante peor que en sus conjuntos de origen, sobre todo
en pieles oscuras y en enfermedades poco comunes, y el ajuste fino con imágenes de fototipos diversos cerró la
brecha~\cite{daneshjou2022disparities}.

\begingroup La disparidad no desaparece cuando el modelo asiste a un médico en lugar de decidir solo. En un
experimento con 389 dermatólogos y 459 médicos de atención primaria, el apoyo de un sistema de aprendizaje
profundo mejoró la exactitud de los dos grupos, pero en los médicos de atención primaria amplió la diferencia
entre pieles claras y oscuras~\cite{groh2024learningaided}.\clubpenalty=10000 \par\endgroup

La población de destino tampoco es la de los conjuntos de entrenamiento. Una revisión sistemática de 272
estudios sobre detección de cáncer de piel con aprendizaje automático encontró que muchos usaron conjuntos de
Europa y Estados Unidos con fototipos mayormente claros, y que solo dos usaron datos de poblaciones de baja
prevalencia, no derivadas a un especialista; por eso los autores no recomiendan todavía adoptar estos modelos en atención
primaria~\cite{jones2022intelligence}. Para la población argentina existe desde 2023 un conjunto público propio,
creado para evaluar herramientas de este tipo en ella~\cite{riccilara2023dataset}.

Con qué datos y con qué protocolo se evaluó un trabajo decide entonces si su resultado se sostiene fuera de su
conjunto de origen, y ese es el criterio con el que la sección siguiente revisa la literatura.
